% -*- mode:LaTex; mode:visual-line; mode:flyspell; fill-column:75-*-

\chapter{Semantic Mapping and Next-Best-View Planning} \label{chapNBV}

% TODO: Chapter intro — tie back to Chapter~\ref{chapDataset}: given a
% 2D detector, how do we build a 3D semantic disease map and plan views
% to efficiently complete that map? Describe the chapter structure.

\section{Semantic OctoMap Representation} \label{secSemanticOctoMap}

% TODO: Describe your 3D semantic map representation in detail.

\subsection{Volumetric Occupancy with OctoMap} \label{secOctoMapBase}

% TODO: Review the OctoMap framework as a foundation.
% - Octree structure; resolution parameter.
% - Probabilistic occupancy with log-odds updates.
% - Memory efficiency for large environments.
% Cite Hornung et al. (2013).

\subsection{Semantic Extension: Disease Class Labels} \label{secSemanticExtension}

% TODO: Describe how each voxel is extended with a semantic state:
% - Class label (healthy / infected / unknown).
% - Per-class confidence from the 2D detector.
% - Bayesian log-odds fusion rule for combining detections across views.
% Include equations for the update rule.
%
% p(c | z_{1:t}) using log-odds:
% l_t = l_{t-1} + log( p(c|z_t) / (1 - p(c|z_t)) ) - log( p(c) / (1-p(c)) )

% \begin{figure}[h]
%   \centering
%   \includegraphics[width=\textwidth]{figs/nbv/semantic_octomap}
%   \caption{TODO: Illustration of the semantic OctoMap: 3D voxel grid
%            colored by disease class confidence.}
%   \label{figSemanticOctoMap}
% \end{figure}

\subsection{Projection from 2D Detections to 3D Voxels} \label{secProjection}

% TODO: Explain how bounding-box or mask detections in image space are
% back-projected into 3D using depth/stereo and the known camera pose.
% - Depth association: stereo disparity or ZED2i depth.
% - Ray casting from camera origin through detected pixels.
% - Handling of detection confidence in the update step.

\section{Simulation Environment} \label{secSimulation}

% TODO: Describe the simulation used to develop and evaluate NBV planners
% before field deployment.

\subsection{Orchard Simulation Setup} \label{secOrchardSim}

% TODO: Describe the Gazebo / PyBullet simulation:
% - Orchard model (tree geometry, row spacing).
% - Simulated disease placement (location, density).
% - Robot and sensor model (camera FOV, noise).

% \begin{figure}[h]
%   \centering
%   \includegraphics[width=0.85\textwidth]{figs/nbv/simulation}
%   \caption{TODO: Screenshot of orchard simulation with robot
%            and simulated fire blight lesions.}
%   \label{figSimulation}
% \end{figure}

\subsection{Ground Truth Generation with Gaussian Splatting} \label{secGSSplat}

% TODO: Describe how Gaussian splatting (3DGS) was used on real data
% to generate high-fidelity 3D ground-truth disease maps for quantitative
% evaluation of the mapping pipeline.
% - Data collection procedure for splat training.
% - Rendering pipeline for ground-truth labels.
% - How ground truth is used to compute mapping accuracy metrics.

\section{Next-Best-View Planning} \label{secNBVPlanning}

% TODO: Describe the NBV planning problem formulation and the three
% planner variants you evaluate.

\subsection{Problem Formulation} \label{secNBVFormulation}

% TODO: Formally define the NBV planning problem.
% - State: current semantic OctoMap, robot pose.
% - Action space: candidate viewpoints (position + orientation).
% - Objective: maximize information gain about disease distribution
%   subject to kinematic and safety constraints.
% - Evaluation metric: map completeness, semantic accuracy, path length.

\subsection{Baseline: Grid Scan Planner} \label{secGridScan}

% TODO: Describe the row-by-row grid scan baseline.
% - Fixed waypoint spacing along orchard rows.
% - No feedback from the current map state.
% - Represents the common manual scouting approach.

\subsection{Volumetric NBV Planner} \label{secVolNBVPlanner}

% TODO: Describe the volumetric information-gain NBV planner.
% - Candidate view sampling strategy.
% - Information gain metric: expected number of unknown voxels observed.
% - Ray-casting implementation for gain estimation.
% - View selection and path planning to next view.

% \begin{figure}[h]
%   \centering
%   \includegraphics[width=\textwidth]{figs/nbv/vol_nbv_diagram}
%   \caption{TODO: Diagram of volumetric NBV planner: candidate views,
%            ray-casting, and gain scores.}
%   \label{figVolNBV}
% \end{figure}

\subsection{Semantic NBV Planner} \label{secSemanticNBVPlanner}

% TODO: Describe your proposed semantic NBV planner.
% - Motivation: prioritize views that reduce uncertainty about disease
%   specifically, not just any unknown space.
% - Semantic information gain: entropy over disease class distribution
%   in voxels visible from candidate view.
% - Weighting: balance occupancy uncertainty and semantic uncertainty.
% - How this differs from volumetric NBV in practice.

% \begin{figure}[h]
%   \centering
%   \includegraphics[width=\textwidth]{figs/nbv/semantic_nbv_diagram}
%   \caption{TODO: Diagram of semantic NBV planner highlighting
%            disease-uncertain regions driving view selection.}
%   \label{figSemanticNBV}
% \end{figure}

\section{Simulation Experiments} \label{secSimExperiments}

% TODO: Describe experiments comparing the three planners in simulation.

\subsection{Experimental Setup} \label{secSimSetup}

% TODO: Describe simulation scenarios:
% - Disease density levels tested.
% - Number of trials per condition.
% - Budget: maximum number of views or total path length.
% - Metrics: map completeness (F1 over voxels), semantic accuracy,
%   detection recall, path efficiency.

\subsection{Results} \label{secSimResults}

% TODO: Report quantitative comparison of Grid Scan vs. Volumetric NBV
% vs. Semantic NBV.
% Include tables and plots (map completeness vs. views, path length).

% \begin{table}[h]
%   \centering
%   \caption{TODO: Planner comparison in simulation.}
%   \label{tabNBVSimResults}
%   \begin{tabular}{lccc}
%     \hline
%     Planner & Map F1 & Semantic Accuracy & Path Length (m) \\
%     \hline
%     Grid Scan & & & \\
%     Volumetric NBV & & & \\
%     Semantic NBV (ours) & & & \\
%     \hline
%   \end{tabular}
% \end{table}

% \begin{figure}[h]
%   \centering
%   \includegraphics[width=\textwidth]{figs/nbv/sim_results}
%   \caption{TODO: Map completeness vs. number of views for all planners.}
%   \label{figSimResults}
% \end{figure}

\subsection{Qualitative Analysis} \label{secSimQualitative}

% TODO: Show example semantic maps produced by each planner at
% comparable view budgets. Discuss where semantic NBV improves over
% volumetric NBV (regions of interest covered earlier).

\section{Summary} \label{secNBVSummary}

% TODO: Summarize:
% - The semantic OctoMap representation and update rule.
% - The three NBV planners and their design rationale.
% - Key simulation result: semantic NBV achieves X\% higher disease
%   recall at the same view budget compared to grid scan and volumetric NBV.
% - These planners are deployed on the real robot in Chapter~\ref{chapErwin}.
